\begin{table}[tp]
\centering
\caption{Stochastic cohypomonotone inclusions: \href{https://github.com/k3nfalt/longtails/blob/main/LongTails/Setups/StochasticCohypomonotoneInclusion.lean}{StochasticCohypomonotoneInclusion}.}\label{tab:StochasticCohypomonotoneInclusion}
\parbox{\linewidth}{\footnotesize In a finite-dimensional real inner-product space, $F$ is $L$-Lipschitz, $G$ is monotone with an everywhere-defined exact resolvent $J_{\gamma G}$ for every $\gamma>0$, and $H=F+G$ satisfies $\langle v-w,x-y\rangle\ge-\rho\|v-w\|^2$ for $v\in H(x),w\in H(y)$. Assume $L>0$, $0\le\rho<1/(2L)$, $R,\sigma\ge0$, $\varepsilon>0$, and $0\in H(x_\star)$ with $\|x_0-x_\star\|\le R$. The jointly measurable oracle $g$ has integrable samples and squared errors, $\int g(x,\omega)\,dp(\omega)=F(x)$ and $\int\|g(x,\omega)-F(x)\|^2\,dp(\omega)\le\sigma^2$. Samples are IID with marginal $p$ under probability law $P$. Target: measurable $(\bar x,u)$ with $u\in G(\bar x)$ on every path and integrable residual satisfying $\int\|F(\bar x)+u\|\,dP\le\varepsilon$.}\par\smallskip
\begin{threeparttable}
\footnotesize
\setlength{\tabcolsep}{4pt}\renewcommand{\arraystretch}{1.25}
\begin{tabularx}{\linewidth}{@{}>{\raggedright\arraybackslash}p{0.20\linewidth}>{\raggedright\arraybackslash}p{0.32\linewidth}>{\raggedright\arraybackslash}p{0.25\linewidth}>{\raggedright\arraybackslash}X@{}}
\toprule
\textbf{Result} & \textbf{Oracle complexity} & \textbf{Novelty} & \textbf{Limitations} \\
\midrule
\textbf{(1)} \href{https://github.com/k3nfalt/longtails/blob/main/LongTails/Theorems/RegularizedProxExtragradientCohypomonotone.lean}{Regularized\allowbreak Prox\allowbreak Extragradient\allowbreak Cohypomonotone} \textbf{(Ours)} $(\star\star\star)$ \href{https://github.com/k3nfalt/longtails/blob/main/scripts/reviews/RegularizedProxExtragradientCohypomonotone.md}{(review)}
& $\mathcal{O}\!\left(\tau^{-8}\left(1+\frac{LR}{\varepsilon}+\frac{\sigma^2}{\varepsilon^2}\right)\ell_\tau^4\right)$ total samples and resolvent calls; fixed budget. Tuning uses $L,\rho,R,\sigma,\varepsilon$.
& & \\
\midrule[1.2pt]
\textbf{(2)} \citep{akw2026v3adjoint}, Remark C.7; stochastic composite case
& $\mathcal{O}\!\left(\frac{L^2R^2\sigma^2}{\delta^4\varepsilon^4}\,Q_2\right)$ total calls, for $LR,\sigma\ge\varepsilon$; $\rho<1/L$. Known $\sigma,\varepsilon$; resolvent-residual accuracy.
& \textbf{(1)}: $\varepsilon^{-2}$ instead of $\varepsilon^{-4}$ at fixed positive parameters, up to logs. For $LR,\sigma\ge\varepsilon$, speedup scales as $\tau^8\min\{LR\sigma^2/\varepsilon^3,L^2R^2/\varepsilon^2\}/\ell_\tau^4$, ignoring \textbf{(2)}\textquotesingle s logs. Strong certificate directly; converting \textbf{(2)} adds $\mathcal{O}(1+\sigma^2/\varepsilon^2)$ samples and one resolvent.
& \textbf{(1)}: $\rho<1/(2L)$ instead of $\rho<1/L$; $\tau^{-8}$ diverges as $\rho\uparrow1/(2L)$ while $\delta^{-4}$ stays bounded. \\
\midrule
\textbf{(3)} RAIN++ \citep{sciChenLuo2024v4}, Theorem 5.4 and Proposition 5.3
& $\mathcal{O}\!\left(\left(1+\frac{LR}{\varepsilon}+\frac{\sigma^2}{\varepsilon^2}\right)Q_3\right)$ expected samples. Special case: $G=0$, $F$ a smooth saddle-gradient operator; $\rho\le1/(2L)$.
& \textbf{(1)}: Scope extends to composite inclusions. Matches polynomial dependence on $L,R,\sigma,\varepsilon$ on the overlap; fixed instead of expected oracle budget.
& \textbf{(1)}: excludes $\rho=1/(2L)$; polynomial $\tau^{-8}$ loss on the overlap. \\
\midrule
\textbf{(4)} AR \citep{sciJiLanZhu2026v2}, Theorem 2.2
& $\mathcal{O}\!\left(1+\frac{LR}{\varepsilon}h+\frac{\sigma^2}{\varepsilon^2}\ell^3h^2\right)$ samples and projections. Special case: monotone $F$, $G=N_X$, closed convex $X$, $x_0\in X$, $\varepsilon\le LR$.
& \textbf{(1)}: Scope extends to nonmonotone cohypomonotone inclusions and general resolvents. Matches polynomial parameter dependence at $\rho=0$.
& \textbf{(1)} at $\rho=0$: $\ell^4$ instead of $h$ on $LR/\varepsilon$, and $\ell^4$ instead of $\ell^3h^2$ on noise; expected norm instead of RMS residual. \\
\midrule
\textbf{(5)} \citep{akw2026v3adjoint}, Corollary 2.2; deterministic composite case
& $\mathcal{O}\!\left(\left(1+\frac{LR}{\delta^2\varepsilon}\right)\left[1+\log\!\left(1+\frac{LR}{\delta\varepsilon}\right)\right]\right)$ total calls. Special case: $\sigma=0$; $\rho<1/L$.
& \textbf{(1)}: extends to stochastic samples; no established rate improvement at $\sigma=0$.
& \textbf{(1)}: narrower $\rho$ range; $\tau^{-8}\ell_\tau^4$ overhead versus \textbf{(5)}'s $\delta^{-2}$ and single accuracy logarithm. \\
\bottomrule
\end{tabularx}
\begin{tablenotes}[flushleft]\footnotesize
\item[] \textbf{Abbr:} $\tau=1-2\rho L$, $\delta=1-\rho L$, $\ell=1+\log(1+LR/\varepsilon)$, $\ell_\tau=\ell+\log(1/\tau)$, $h=1+\log\ell$. $Q_2$ denotes the unspecified polylogarithmic factor in $L,R,\sigma,\varepsilon^{-1},\delta^{-1}$ in Remark C.7; $Q_3=\operatorname{polylog}(1+LR/\varepsilon,1+\varepsilon^{-1},1+R,1+\sigma)$ denotes the stated logarithmic overhead of RAIN++. Polylogarithms are nonvanishing. One call is one $F$ sample or one exact $G$ resolvent; all batches count. $N_X$ is the normal cone, whose resolvent is projection onto $X$; RMS means root mean square. Resolvent residual means $\eta^{-1}\|x-J_{\eta(F+G)}x\|$, with $\eta=(1+\rho L)/(2L)$ in \textbf{(2)}, \textbf{(5)}.
\end{tablenotes}
\end{threeparttable}
\end{table}
\clearpage
